\section{Introduction}

The deployment of autonomous multi-agent systems marks a structural shift in how complex goals are pursued in software infrastructure. Unlike single-agent designs, where a human operator can interrupt a running process at will, multi-agent architectures distribute intent across concurrent, interdependent agents that negotiate, delegate, and act without sustained human oversight. This distribution is a feature — it enables scale, parallelism, and emergent problem-solving that no individual agent can achieve alone. It is also, we argue, the point at which a previously optional property becomes architecturally mandatory: the capacity for structured, accountable withdrawal.

We call this property the \textbf{Bailout Protocol}, and this paper is its definition, formalization, and empirical evaluation. In the BX3 framework --- the conceptual substrate on which this work rests --- every agentic system is characterized by four irreducible constituents: \textbf{purpose}, the goal or objective state toward which the system acts; \textbf{bounds}, the envelope of permissible actions and states that must not be exceeded; \textbf{fact}, the verified, attributable record of events, decisions, and their causal lineage; and the agent itself as the entity that relates purpose to action within bounds, producing fact as its observable output. These four terms are not metaphorical. They are the minimum semantic vocabulary required to reason about what a system is doing, why it is doing it, and whether what it is doing is acceptable.

The central observation of this paper is that purpose, bounds, and fact are insufficient without a fourth mechanism: a \textbf{bail-out} --- a structured, non-negotiable fallback that activates when the relationship between purpose and bounds breaks down. Without such a fallback, the system's fact record either continuesunchecked, producing an accountable-seeming log of an increasingly non-compliant system, or it stalls entirely, leaving no record at all of what went wrong and why the system stopped. Neither outcome satisfies the demands of accountability. An accountable system is not one that records everything; it is one that can be stopped --- and that stopping is guaranteed by architecture, not by hope.

\subsection{Why Multi-Agent Systems Require Mandatory Bail-Out}

Multi-agent systems are particularly exposed to this failure class because their dynamics are emergent. An individual agent may operate within its stated purpose and declared bounds on every step, yet the composition of many such agents acting concurrently can produce aggregate behavior that violates bounds in ways no single-agent reasoning would predict. Cascade effects, resource contention, and distributed decision-making compound unpredictably. In such environments, the risk is not localized --- it propagates across the system in real time.

Existing autonomous systems tend to address this risk through one of three inadequate approaches. The first is \textbf{operational oversight}: a human monitor observes system behavior and intervenes manually. This approach does not scale, introduces latency proportional to human response time, and fundamentally violates the premise of autonomous operation. The second is \textbf{static policy gates}: pre-defined rule sets that halt the system when certain conditions are met. While more scalable than human oversight, static gates are brittle --- they can only act on conditions they were explicitly programmed to anticipate, and they offer no semantics for interpreting whether a triggered halt represents a failure, a success, or a neutral boundary crossing. The third, and most common, is \textbf{implicit trust}: the system is assumed to be sufficiently well-designed that out-of-bounds behavior will not occur. This is the most dangerous approach, because it produces no failure record when the assumption inevitably breaks.

What none of these approaches provide is a \textbf{mandatory architectural fallback} --- a mechanism that is intrinsic to the system's design, independent of any individual agent's compliance, and guaranteed to produce a fact-accurate record of the system's terminal state regardless of how or why the bailout was triggered. This guarantee is what the Bailout Protocol is designed to supply.

\subsection{Failure Modes in the Absence of a Bailout Protocol}

When a multi-agent system operates without a mandatory bail-out mechanism, three distinct failure modes arise, each of which undermines accountability in a different way.

\textbf{Drift without detection.} Purpose and bounds diverge gradually, below the threshold of any static policy gate, until the system's actions are functionally misaligned with its stated intent. The fact record appears coherent --- events are logged, decisions are attributed --- but the aggregate direction is non-compliant. Without a bail-out, there is no architectural signal that distinguishes this state from normal operation. The system is not broken in any locally detectable sense, but it has drifted beyond its authorized purpose.

\textbf{Collapse without record.} The system encounters a boundary condition it cannot resolve --- a resource exhaustion, a cascade failure, an unresolvable inter-agent conflict --- and terminates abruptly. In the absence of a structured bail-out, the terminal state is not captured. The fact record ends mid-event, leaving investigators or auditors with a partial log that cannot be attributed to a specific cause. The system is neither in compliance nor in violation --- it is simply gone.

\textbf{Compliance theater.} The system possesses a bail-out mechanism that is agentically implemented but not architecturally guaranteed --- that is, each agent may choose to respect it, but no structural property forces compliance. In this mode, the fact record shows a bailout having occurred, but it is indistinguishable from any other logged event: there is no way to verify that the bail-out was triggered by the correct condition, that it executed fully, or that it was not suppressed by a compromised agent. The mechanism exists but provides no accountable guarantee.

Each of these failure modes is empirically observable in existing multi-agent deployments, and each is directly traceable to the absence of a mandatory, architecturally enforced bail-out protocol. The Bailout Protocol we propose is specifically designed to close all three gaps.

\subsection{The Core Insight: Accountability Requires Architectural Fallback}

The theoretical contribution of this paper rests on a single, generalizable insight: \textbf{accountability is not a property of agents --- it is a property of architecture}. An agent can act within its purpose and within its declared bounds; it can produce a rich, detailed fact record; and it can do all of this while participating in a system that is nonetheless misaligned at the aggregate level, or that will catastrophically fail when a specific class of boundary condition is encountered. Individual agent compliance does not and cannot guarantee system-level accountability.

This means that accountability must be enforced at a level of abstraction above the individual agent. In the BX3 framework, this enforcement takes the form of a \textbf{bounds substrate} --- a structural layer that is independent of any agent's internal logic, that has the authority to interrupt the purpose-fact loop, and that is guaranteed to fire when a defined class of boundary condition is detected. The Bailout Protocol is the specification of this bounds substrate: its trigger conditions, its execution semantics, its output guarantees, and its relationship to the fact record.

We emphasize that the Bailout Protocol is not an exception handler. It is not a try-catch block or a circuit breaker in the conventional sense. It is a \textbf{constitutive element of the system's identity}: the mechanism without which purpose has no meaning, bounds have no force, and fact has no anchor to reality. A system that lacks a mandatory bail-out is not a malfunctioning instance of an accountable system --- it is a different class of system, one that cannot be held accountable regardless of what its fact record contains.

\subsection{Paper Roadmap}

The remainder of this paper is organized as follows. Section~\ref{sec:formal-model} presents the formal model of the Bailout Protocol within the BX3 framework, defining its constituent terms --- purpose, bounds, fact, and the bail-out trigger --- with mathematical rigor. Section~\ref{sec:specification} provides the architectural specification of the protocol, including trigger taxonomy, execution semantics, and the guaranteed outputs that constitute the fact-accurate terminal record. Section~\ref{sec:implementation} describes the reference implementation and the specific design decisions that preserve the protocol's guarantees in the presence of concurrent agent execution, network partition, and resource exhaustion. Section~\ref{sec:evaluation} presents empirical results from three deployment scenarios that exercise each of the three failure modes described above, demonstrating that the Bailout Protocol correctly detects, halts, and records each case. Section~\ref{sec:related-work} situates this contribution relative to existing literature on agent accountability, multi-agent safety, and interrupt mechanisms in autonomous systems. Section~\ref{sec:conclusion} summarizes the contribution and identifies the open questions that remain.
